\section{Conclusion}
\label{sec:conclusion}

Asking a language model what changed, in a closed vocabulary, and testing the answer on later
data is a sound interface; letting the test switch the answer on is not enough: the registered
gate acted late or not at all, on whatever size the model stated. Certify-then-hedge keeps the
certificate and changes its use: the \eprocess{} prices hypotheses, the data size them, and the
controller hedges at the critical fractile of the resulting mixture. Under a known null and arrival law that
rule keeps the gate's error budget as a bound on expected exposure for any model output, and in
exploratory tests registered before their live runs it beat the gate at all three of the
benchmark's cost ratios and the operations research baseline at the middle and highest, with two
proposer models, though at the highest non-inferiority on no-shock episodes was not
established and on noisy lead times the gain over the baseline was not confirmed. On this
simulator the value came mostly from mechanisms the alert-timed content-free control shares,
concentrated in lead-time shifts; the text cut commitment on no-shock episodes and, when
stockouts dominate, added value with one of the two models. Whether text earns more where it
says what the data reveal only late, under a null the guarantee fully covers and on
InventoryBench's own stochastic lead times and lost shipments, is the question this design can
now answer.
